\documentclass[12pt, a4paper]{article}

% --- Essential Packages ---
\usepackage{amsmath, amssymb, amsfonts} 
\usepackage{geometry}                   
\geometry{hmargin=1.5cm, vmargin=2cm}
\usepackage{hyperref}                   
\hypersetup{
    colorlinks=true,
    linkcolor=blue,
    filecolor=magenta,      
    urlcolor=cyan,
}
\usepackage{enumitem}
\usepackage{ifthen} 
\usepackage{needspace} % <--- ADDED: Package to require a specific amount of space
\usepackage{tabularray}

\newcommand{\N}{\mathbb{N}}
\newcommand{\Z}{\mathbb{Z}}
\newcommand{\Q}{\mathbb{Q}}
\newcommand{\R}{\mathbb{R}}
\newcommand{\C}{\mathbb{C}}
\newcommand{\F}{\mathbb{F}}

% --- Custom Streamlined Solution Environments ---

\newcounter{probnum} 
\setcounter{probnum}{0} 

% 1. Standard solution environment
\newenvironment{solution}[1]{%
    \vspace{1em}% 
    % --- ADDED: Checks if there is enough space. If not, forces a new page. ---
    \needspace{10\baselineskip}% 
    \ifthenelse{\equal{#1}{}}%
        {\stepcounter{probnum}}% 
        {\setcounter{probnum}{#1}}% 
    \begin{list}{}{%
        \setlength{\leftmargin}{2em}% 
        \setlength{\labelwidth}{1.5em}% 
        \setlength{\labelsep}{0.5em}% 
        \setlength{\topsep}{0pt}%
        \setlength{\partopsep}{0pt}%
    }%
    \item[\textbf{\theprobnum.}]%
}{%
    \end{list}%
    \par\vspace{1em}% 
}

% 2. Sub-parts enumeration template
\newlist{subparts}{enumerate}{1}
\setlist[subparts]{
    label=\textbf{\alph*.},
    leftmargin=1.3em, 
    labelwidth=0.8em, 
    labelsep=0.5em,   
    align=left,       
    topsep=0pt,
    itemsep=0.2em
}

% 3. Dedicated environment for solutions 
\newenvironment{solutionparts}[1]{%
    \vspace{1em}%
    % --- ADDED: Checks if there is enough space. If not, forces a new page. ---
    \needspace{10\baselineskip}% 
    \ifthenelse{\equal{#1}{}}%
        {\stepcounter{probnum}}%
        {\setcounter{probnum}{#1}}%
    \noindent\makebox[1.5em][r]{\textbf{\theprobnum.}}\vspace*{-\baselineskip}%
    \begin{subparts}[leftmargin=3.3em, labelwidth=0.8em, labelsep=0.5em, align=left]
}{%
    \end{subparts}%
    \par\vspace{1em}%
}

\begin{document}

% ==========================================
%                TITLE PAGE
% ==========================================
\begin{titlepage}
    \centering
    \vspace*{2in}
    {\Huge \bfseries Solutions Manual \par}
    \vspace{0.5in}
    {\Large Abstract Algebra by Dummit \& Foote (3rd ed.) \par}
    \vspace{1.5in}
    {\large \textbf{Prepared by:} \par}
    {\large Daniel Jiesung Jung (\texttt{daniel72@postech.ac.kr})\par}
    \vfill
    {\large \today \par}
\end{titlepage}

% ==========================================
%              TABLE OF CONTENTS
% ==========================================
\newpage
\tableofcontents
\newpage

% ==========================================
%                PREFACE
% ==========================================
\newpage
\section*{Preface}
\addcontentsline{toc}{section}{Preface} % Adds the Preface to the Table of Contents

This solutions manual grew out of my effort to study Dummit \& Foote thoroughly.

The textbook is famous for placing important results in the exercises rather than in the main body of the text. Because of this, I will summarize any theory-like results at the beginning of each subsection's solutions so they can be referenced easily.

Please reach out to me at \texttt{daniel72@postech.ac.kr} if you spot any mistakes or have any suggestions for improvement.


\vspace{1cm}
\begin{flushright}
Daniel Jiesung Jung \\
Pohang, South Korea \\
\today
\end{flushright}
\newpage

% ==========================================
%                MAIN CONTENT
% ==========================================


% ==========================================
%                 Section 1
% ==========================================
\section{Introduction to Groups} 
\subsection{Basic Axioms and Examples}

Skipped problems: 1 \(\sim\) 9, 11 \(\sim\) 14

\noindent Problems that need revisiting: 19
\vspace{1em}


%Recall the residue classes in \(\Z / n\Z\).
%\begin{itemize}
%    \item The elements are the equivalence classes of \(\sim\) where \(a \sim b\) if and only if \(n \mid b-a\).
%          There are exactly \(n\) elements in \(\Z / n\Z\): \(\{\bar 0, \bar 1, \cdots, \overline{n-1}\}\).
%    \item Addition is defined as \(\overline{a} + \overline{b} = \overline{a+b}\).
%    \item Multiplication is defined as \(\overline{a} \cdot \overline{b} = \overline{ab}\).
%    \item These binary operations are well-defined.
%\end{itemize}

%\begin{solution}{3}
%    \[
%    \begin{aligned}
%        (\overline{a} + \overline{b}) + \overline{c} &= \overline{a+b} + \overline{c}\\
%                                                     &= \overline{(a+b)+c}\\
%                                                     &= \overline{a+(b+c)} & (\because \text{associativity in } \Z)\\
%                                                     &= \overline{a} + \overline{b+c} \\
%                                                     &= \overline{a} + (\overline{b} + \overline{c})
%    \end{aligned}
%    \]
%\end{solution}

\begin{solution}{10}
    Let \(G = \{g_1, \cdots, g_n\}\) be a finite group and \(A\) the group table of \(G\).
    By definition, \((A)_{ij} = g_ig_j\). 

    If \(G\) were abelian, then for any indices \(i, j\), we have \(g_ig_j = g_jg_i\), which implies that
    \((A)_{ij} = (A)_{ji}\). The choices for \(i\) and \(j\) were arbitrary, so it follows that \(A\) is symmetric.

    Conversely, if \(A\) were symmetric then for any indices \(i, j\) we have \((A)_{ij} = (A)_{ji}\),
    which in turn implies that \(g_ig_j = g_jg_i\), concluding that \(G\) is abelian. \(\blacksquare\)
\end{solution}


\begin{solution}{15}
    We do this by induction.
    
    Base case: If \(n=1\), then it is evident that \((a_1)^{-1} = a_1^{-1}\).
    
    Inductive step: If \((a_1a_2\cdots a_k)^{-1} = a_k^{-1}a_{k-1}^{-1}\cdots a_1^{-1}\) for some integer \(k\), then we see that
    \[
    \begin{aligned}
        \left((a_1a_2\cdots a_k) a_{k+1}\right)^{-1}  &= a_{k+1}^{-1} (a_1a_2\cdots a_k)^{-1}\\
                                                      &= a_{k+1}^{-1} (a_k^{-1} a_{k-1}^{-1}\cdots a_1^{-1})\\
                                                      &= a_{k+1}^{-1} a_k^{-1} a_{k-1}^{-1}\cdots a_1^{-1},
    \end{aligned}
    \]
    which completes the proof. \(\blacksquare\)
\end{solution}

\begin{solution}{16}
    If \(x^2 = 1\), then \(x\) has order of at most 2 by definition. Conversely, if \(x\) has order 1 then \(x^2 = 1^2 = 1\), and if \(x\) has order 2 then \(x^2 = 1\). \(\blacksquare\)
\end{solution}

\begin{solution}{17}
    If \(|x|=n\), then we have \(x^n = 1\). Multiplying by \(x^{-1}\) on both sides, we get \(x^{n-1} = x^{-1}\). \(\blacksquare\)
\end{solution}

\begin{solution}{18}
    For the forward directions,
    from \(xy = yx\) multiply \(y^{-1}\) on the left of both sides to deduce \(y^{-1}xy = x\),
    then multiply \(x^{-1}\) on the left of both sides to deduce \(x^{-1}y^{-1}xy = 1\).
    For the backward directions,
    from \(x^{-1}y^{-1}xy = 1\) multiply \(x\) then \(y\) on the left of both sides to retrieve
    \(y^{-1}xy = x\) and then \(xy = yx\). \(\blacksquare\)
\end{solution}

\begin{solutionparts}{19}
    \item This can be done by simple counting.
    \[
    \begin{aligned}
        x^{a+b} &= \underbrace{x\cdots\cdots \cdots x}_{a+b\text{ times}} \\
                &= \underbrace{x \cdots x}_{a\text{ times}} \underbrace{x \cdots x}_{b\text{ times}}\\
                &= (\underbrace{x \cdots x}_{a\text{ times}}) (\underbrace{x \cdots x}_{b\text{ times}})\\
                &= x^ax^b.
    \end{aligned}
    \]
    also,
    \[
    \begin{aligned}
        (x^a)^b &= (\underbrace{x\cdots x}_{a\text{ times}})^b\\
                &= \underbrace{(\underbrace{x\cdots x}_{a\text{ times}})\cdots(\underbrace{x\cdots x}_{a\text{ times}})}_{b\text{ times}}\\
                &= \underbrace{x\cdots \cdots x}_{ab\text{ times}}\\
                &= x^{ab}.
    \end{aligned}
    \]

    \item This can be seen by simply observing that
    \(x^{-a}x^a = \underbrace{x^{-1}\cdots x^{-1}}_{a\text{ times}}\underbrace{x \cdots x}_{a \text{ times}} = 1\), and \(x^a x^{-a}\) similarly.

    \item This seems to be way too much case classification without much reward, so I will skip this at the moment.
\end{solutionparts}

\begin{solution}{20}
    
    For any positive integer \(k\), we have \(x^k = 1\) if and only if \(1 = (x^{-1})^k\),
    by multiplying appropriate terms in both sides.
    Now, there are 2 cases: \(x\) can have either finite or infinite order.

    If \(x\) has finite order \(n\), then we have \(x^k \neq 1\) for \(k = 1, \cdots, n-1\),
    and this implies that \(1 \neq (x^{-1})^k\).
    Also, we have \(x^n = 1\), which implies \(1 = (x^{-1})^n\).
    This shows that \(x^{-1}\) has order \(n\).
    

    Now, if \(x\) has infinite order, then for any fixed positive integer \(n\) we have \(x^n \neq 1\).
    This implies that \(1 \neq (x^{-1})^n\),
    which implies that \(x^{-1}\) has infinite order.

    This completes the proof. \(\blacksquare\)
\end{solution}

\begin{solution}{21}
    If \(x\) has odd order in the group, then we have \(x^n = 1\) for some odd integer \(n\).
    Then, \(x^{n+1} = x\).
    Since \(n+1\) is even, we can write \(n+1 = 2k\) for some integer \(k\).
    From this we conclude that \(x^{2k} = (x^2)^k= x\), completing the proof. \(\blacksquare\)
\end{solution}

\begin{solution}{22}
    We first observe that for any positive integer \(i\), we have \((g^{-1}xg)^i = g^{-1}x^ig\). (Technically this must be proven through induction, but we will not do it here.)

    If \(x\) has finite order \(n\), then we have \(x^k \neq 1\) for \(k = 1, \cdots, n-1\), and \(x^n = 1\). For any \(k \in \{1, \cdots, n-1\}\), we have \((g^{-1}xg)^k = g^{-1}x^kg\).
    If this were 1, then it would imply that \(x^k = gg^{-1} = 1\), which must not be true,
    so we have \(g^{-1}x^kg \neq 1\). Also, we observe that \((g^{-1}xg)^n = g^{-1}x^ng = g^{-1}g = 1\).
    This shows that \(g^{-1}xg\) also has order \(n\). If \(x\) has infinite order, then we can prove similarly that \(g^{-1}xg\) has infinite order too.

    To conclude \(|ab| = |ba|\) from this, simply let \(x = ab\) and \(g = a\). \(\blacksquare\)
\end{solution}

\begin{solution}{23}
    Observe that \((x^s)^t = x^{st} = x^n\). Observe further two things.
    \begin{itemize}
        \item \((x^s)^k \neq 1\) for \(k = 1, \cdots, t-1\), since here \(sk < st = n\),
        and \(x^i \neq 1\) for all \(1 \le i < n\). 
        \item \((x^s)^t = 1\), since \((x^s)^t = x^n\).
    \end{itemize}
    This concludes that \(x^s\) has order \(t\). \(\blacksquare\)
\end{solution}

\begin{solution}{24}
    Pick any two elements \(a, b \in G\) that commutes.
    We first handle the obvious case when \(n = 0\): \((ab)^0 = a^0b^0 = 1\).
    
    Following the hint, we will first do induction on positive \(n\).
    For the base case, if \(n=1\) then \((ab)^1 = a^1b^1\) is trivially true.
    For the inductive step, assume that \((ab)^k = a^kb^k\).
    Then, we have
    \[
    \begin{aligned}
        (ab)^{k+1} &= a(ba)^kb\\
                   &= a(ab)^kb &\text{ (since \(a\) and \(b\) commute)}\\
                   &= aa^kb^kb &\text{ (by the induction hypothesis)}\\
                   &= a^{k+1}b^{k+1}.
    \end{aligned}
    \]
    This shows that for any positive \(n\), we have \((ab)^n = a^nb^n\).

    To show that this fact extends to negative integers,
    fix an arbitrary positive integer \(m\), and we wish to show that \((ab)^{-m} = a^{-m}b^{-m}\).
    We first see that since \(a\) and \(b\) commutes, \(a^{-1}\) and \(b^{-1}\) also commutes, since
    \(a^{-1}b^{-1} = (ba)^{-1} = (ab)^{-1} = b^{-1}a^{-1}\). Then, we observe that
    \[
    \begin{aligned}
        (ab)^{-m}&= ((ab)^{-1})^m\\
                 &= (b^{-1}a^{-1})^m\\
                 &= (a^{-1}b^{-1})^m & \text{ (since \(a^{-1}\) and \(b^{-1}\) commutes)}\\
                 &= (a^{-1})^m(b^{-1})^m & \text{ (again, since \(a^{-1}\) and \(b^{-1}\) commutes)}\\
                 &= a^{-m}b^{-m}.
    \end{aligned}
    \]
    This completes the argument. \(\blacksquare\)
\end{solution}

\begin{solution}{25}
    \(x^2 = 1\) implies that \(x = x^{-1}\). Then, for any \(a, b \in G\), we see that 
    \[
    \begin{aligned}
        ab &= a^{-1}b^{-1}  & \text{ (since \(x = x^{-1}\))}\\
           &= (ba)^{-1}\\
           &= ba.            & \text{ (since \(x = x^{-1}\))}
    \end{aligned}
    \]
    This shows that \(G\) is abelian. \(\blacksquare\)
\end{solution}

\begin{solution}{26}
    We first see that \(\star\) is indeed a binary operation on \(H\),
    i.e. a function such that \(\star: H \times H \to H\),
    since \(H\) is closed under \(\star\).
    (To spell it out, it is because for any \(h, k \in H\), we have \(h \star k \in H\).)

    Now we check the group axioms. Associativity of \(\star\) in \(H\) is inherited directly from the fact that \(G\) is associative.
    If \(1\) is the identity of \(G\), then in particular, for any \(h \in H\), we have \(h\star 1 = 1 \star h = h\).
    Additionally we have to check that \(1 \in H\). To see this, for some \(h \in H\) we have \(h^{-1} \in H\),
    so consequently \(1 = hh^{-1} \in H\).
    Finally, since \(H\) is closed under inverses, for any \(h \in H\) we have \(h^{-1} \in H\), so the inverse axiom is satisfied. This completes the proof. \(\blacksquare\)
\end{solution}

\begin{solution}{27}
    Denote \( \langle x \rangle = \{x^n \mid n \in \Z\}\).
    It suffices to show that for any two elements \(p, q \in \langle x \rangle\), it must be that \(pq \in \langle x \rangle\) and \(p^{-1} \in \langle x \rangle\).
    
    Noting that \(p = x^{n_p}\) for some \(n_p \in \Z\) and \(q = x^{n_q}\) for some \(n_q \in \Z\),
    it is easy to see that \(pq = x^{n_p}x^{n_q} = x^{n_p + n_q} \in \langle x \rangle\) and
    \(p^{-1} = (x^{n_p})^{-1} = x^{-n_p} \in \langle x \rangle\). \(\blacksquare\)
\end{solution}

\begin{solutionparts}{28}
    \item We have
        \[[(a_1, b_1)(a_2, b_2)](a_3, b_3) = ([a_1a_2]a_3, [b_1b_2]b_3)\]
    and
        \[(a_1, b_1)[(a_2, b_2)(a_3, b_3)] = (a_1[a_2a_3], b_1[b_2b_3]).\]
    The associativity of groups \(A\) and \(B\) imply that the two expressions the same, so that \(A \times B\) is associative.

    \item For any \(a \in A\) and \(b \in B\), we have \((1,1)(a,b) = (1a, 1b) = (a,b)\) and \((a,b)(1,1)=(a1, b1) = (a,b)\), implying that \((1,1)\) is the identity of \(A \times B\).
    
    \item For any \(a \in A\) and \(b \in B\), if \(a^{-1}\) is the inverse of \(a\) in \(A\) and \(b^{-1}\) the inverse of \(b\) in \(B\), then we see that \((a,b)(a^{-1}, b^{-1}) = (aa^{-1}, bb^{-1}) = (1,1)\) and \((a^{-1}, b^{-1})(a, b) = (a^{-1}a, b^{-1}b) = (1, 1)\), showing that \((a^{-1}, b^{-1})\) is the inverse of \((a,b)\) in \(A \times B\). \(\blacksquare\)
\end{solutionparts}

\begin{solution}{29}
\(A \times B\) is an abelian group if and only if for any elements \(a_1, a_2 \in A\) and \(b_1, b_2 \in B\), we have \((a_1, b_1)(a_2, b_2) = (a_2, b_2)(a_1, b_1)\). Here, \((a_1, b_1)(a_2, b_2) = (a_1a_2, b_1b_2)\) and \((a_2, b_2)(a_1, b_1) = (a_2a_1, b_2b_1)\). This is true if and only if \(a_1a_2 = a_2a_1\) and \(b_1b_2 = b_2b_1\), if and only if \(A\) and \(B\) are both abelian. \(\blacksquare\)
\end{solution}

\begin{solution}{30}
    We see that \((a, 1)(1, b)\) and \((1, b)(a, 1)\) both evaluate to \((a, b)\), so they commute.
    To see that the order of \((a, b)\) is the lcm of \(|a|\) and \(|b|\), we observe that
    \[
    \begin{aligned}
        (a, b)^n &= [(a,1)(1,b)]^n\\
                 &= (a,1)^n(1,b)^n\\
                 &= (a^n, 1)(1, b^n)\\
                 &= (a^n, b^n),
    \end{aligned}
    \]
    since \((a,1)\) and \((1, b)\) commute.
    For this to be the identity we see that both \(a^n\) and \(b^n\) must be the identity.
    \(a^n\) is the identity if and only if \(|a|\) divides \(n\), and
    \(b^n\) is the identity if and only if \(|b|\) divides \(n\).
    Now, the smallest positive \(n\) such that both \(|a|\) and \(|b|\) divides \(n\),
    must be the lcm of \(|a|\) and \(|b|\),
    and therefore \(|(a, b)|\) must be the lcm of \(|a|\) and \(|b|\). \(\blacksquare\) 
\end{solution}

\begin{solution}{31}
    Let \(G\) be a group of even order.
    We have to prove that there is an element in \(G\) of order 2.

    Following the hint, define the set \(t(G) = \{g \in G \mid g \neq g^{-1}\}\).
    We first see that if some \(a\) is some element in \(G\) such that \(a \in t(G)\),
    then it must be that \(a^{-1} \in t(G)\) since \((a^{-1})^{-1} \neq a^{-1}\).
    Moreover, \(a\) and \(a^{-1}\) must be distinct, again because of the condition \(a \neq a^{-1}\).
    From this we can argue that \(t(G)\) as a set has even order, since the elements in \(t(G)\)
    can be paired in two: \((a, a^{-1})\).
    Since \(G\) also has even order, we conclude that \(G \smallsetminus t(G)\) must also have even order.


    Now, we see that \(x \in G \smallsetminus t(G)\) if and only if \(x = x^{-1}\).
    This implies that \(x^2 = 1\), so if \(x\) is not the identity then \(x\) has order 2 in \(G\).
    Also, since \(1 = 1^{-1}\), we have \(1 \in G \smallsetminus t(G)\).
    Since \(G \smallsetminus t(G)\) has even order,
    some other element in \(G \smallsetminus t(G)\), other than the identity, must exist,
    and this element must have order 2 in \(G\).
    This concludes the proof. \(\blacksquare\)
\end{solution}

\begin{solution}{32}
    By definition of an order of an element, if \(n\) is the order of the element \(x\), we have \(x^k \neq 1\) for \(k = 1, \cdots, n-1\), and \(x^n = 1\). Now, assume for the sake of contradiction that \(x^i = x^j\) for some \(1 \le i < j \le n-1\). Then, \(x^i x^{n-j} = x^jx^{n-j}\) implies \(x^{n+i-j} = x^n\). Observing that \(0 < n+i-j < n\) and \(x^n = 1\), we see that \(x^{n+i-j} = 1\), which contradicts our very first observation. Therefore, \(1, x, \cdots, x^{n-1}\) must be distinct elements.

    To deduce \(|x| \leq |G|\) from here, we simply observe that \(G\) has at least \(n\) elements from the previous paragraph. \(\blacksquare\)

\end{solution}

\begin{solution}{33}
    Assume \(x\) is of order \(n\) in the group \(G\).
    We have \(x^n = 1\) and \(x^{k} \neq 1\) for \(k = 1, \cdots, n-1\).
    Then, for \(1 \le j < 2n\),
    the only value of \(j\) that satisfies \(x^j = 1\)
    is \(j = n\).
    \begin{subparts}
        \item If \(n\) is odd, then for \(1 \le j < 2n\) the only value that satisfies \(x^j = 1\)
        is \(n\), which is odd.

        Now, suppose for the sake of contradiction that
        for some \(i\) such that \(1 \le i < n\), we have \(x^i = x^{-i}\).
        Then, we have \(x^{2i} = 1\). But \(1 \le 2i < 2n\) and \(2i\) is even,
        contradicting our initial observation.
        Therefore, if \(n\) is odd then it must be that
        \(x^i \neq x^{-i}\) for \(i = 1, \cdots, n-1\).

        \item If \(n\) is even, then for \(1 \le j < 2n\) the only value that satisfies \(x^j = 1\) is \(n\),
        and if this is the case then \(j = n = 2k\).
        
        Here, suppose that
        for some \(i\) such that \(1 \le i < n\), we have \(x^i = x^{-i}\).
        Then, we have \(x^{2i} = 1\). But since \(1 \le 2i < 2n\),
        it must be the case that \(2i = 2k\), or that \(i = k\).
        Conversely, we obviously have \(x^k = x^{-k}\) since \(x^{2k} = x^n = 1\).
    \end{subparts}
    This completes the argument. \(\blacksquare\)
\end{solution}

\begin{solution}{34}
    \(x\) being an infinite order element in \(G\) implies that
    for any positive integer \(n\), we have \(x^n \neq 1\).
    By similar logic from the solution for problem 32,
    for the sake of contradiction assume that \(x^i = x^j\) for some distinct integers \(i\) and \(j\),
    and assume \(i < j\) without loss of generality.
    Then, \(x^ix^{-i} = x^jx^{-i}\) implies that \(x^{j-i} = 1\), which contradicts our very first observation.
    Therefore, elements of the form \(x^n\) where \(n \in \mathbb Z\) must be distinct. \(\blacksquare\)
\end{solution}

\begin{solution}{35}
    By the definition of the order of an element in a group,
    if \(x\) has finite order \(n\) in \(G\) then we have distinct elements
    \(1, x, \dots, x^{n-1} \in G\), and then \(x^n = 1\).

    Now, fix any \(a \in \Z\), and we are left to show that \(x^a \in \{1, x, \cdots, x^{n-1}\}\).
    By the division algorithm, we can write \(a = qn + r\) for some unique \(q\) and \(r\) where \(0 \le r < n\). Then we observe that 
    \[
    \begin{aligned}
        x^a &= x^{qn + r}\\
            &= (x^n)^q x^r \\
            &= x^r \in \{1, x, \cdots, x^{n-1}\},
    \end{aligned}
    \]
    since \(0 \le r < n\). This completes our argument. \(\blacksquare\)
\end{solution}

\begin{solution}{36}
    The contraposition of the cancellation laws state that
    for \(u, v \in G\), if \(u \neq v\) then we have \(au \neq av\) and \(ua \neq va\) for any \(a \in G\).
    This implies that the group table must not have duplicate elements in the same row or column.
    Keeping this rule in mind, and ensuring that there are no elements of order 4,
    we see that the table below is the only possible group table for \(G = \{1, a, b, c\}\).

    \begin{center}
    \begin{tblr}{
        rows = {1cm, rowsep = 2pt},
        columns = {1cm, colsep = 2pt},
        cells = {m, c},
        hlines, vlines,
        hline{2} = {1.5pt},
        vline{2} = {1.5pt}
    }
    \(\cdot\) & \(1\) & \(a\) & \(b\) & \(c\) \\
    \(1\)     & \(1\) & \(a\) & \(b\) & \(c\) \\
    \(a\)     & \(a\) & \(1\) & \(c\) & \(b\) \\
    \(b\)     & \(b\) & \(c\) & \(1\) & \(a\) \\
    \(c\)     & \(c\) & \(b\) & \(a\) & \(1\)  
    \end{tblr}
    \end{center}
    This matrix is symmetric, so by Exercise 1.1.10 \(G\) must be abelian. \(\blacksquare\)

    (Note: we haven't proved that the group axioms for the binary operation hold for that ``group table''.
    We can see that the identity and inverse axioms hold, but associativity is trickier to see.
    However, the group \(\Z / 2\Z \times \Z / 2\Z\) is isomorphic to this table,
    and we can deduce associativity from there.)
\end{solution}




\end{document}
